\chapter{微分学}\label{ch:b1:derivative}

导数在高中分册中已经计算过多次；缺少的是
一串把计算转化为函数信息的定理，
即罗尔定理、中值定理，
以及它们的推论——单调性判别法、利普希茨 界、
凸性。本章所有内容均针对
定义在\omterm{prop:b1:reals:intervals}{区间} $I$ 上的函数。

\section{导数}

\begin{definition}\label{def:b1:derivative:def}
称 $f \colon I \to \R$ 在 $x_0 \in
I$ 处\emph{可导}\index{可导}，当差商
$\frac{f(x) - f(x_0)}{x - x_0}$ 在 $x \to x_0$ 时存在（有限）极限；
该极限记为 $f'(x_0)$。等价地：
\[
f(x_0 + h) = f(x_0) + f'(x_0)\,h + h\,\varepsilon(h),
\qquad \varepsilon(h) \xrightarrow[h \to 0]{} 0 ,
\]
此时函数图像存在切线 $y = f(x_0) + f'(x_0)(x -
x_0)$。在 $x_0$ 处可导蕴含在 $x_0$ 处\omterm{def:b1:continuity:continuous}{连续}（由上式即可看出）。
$f$ 在 $I$ 上可导，是指它在每一点都可导；
称 $f$ 为 $C^1$ 类函数，若又 $f'$ \omterm{def:b1:continuity:continuous}{连续}；
$C^k$ 类函数\index{$C^k$ 类函数}是指 $f$ 可作 $k$ 次求导，
且 $f^{(k)}$ \omterm{def:b1:continuity:continuous}{连续}。
\end{definition}

\begin{example}\label{ex:b1:derivative:notc1}
“\omterm{def:b1:derivative:def}{可导} $\Rightarrow$ \omterm{def:b1:continuity:continuous}{连续}”的逆命题不成立：
反例为 $\abs{\,\cdot\,}$ 在 $0$ 处。更令人惊讶的是，\omterm{def:b1:derivative:def}{可导}
并不蕴含 $C^1$：函数 $f(x) = x^2 \sin\frac 1x$（$f(0) = 0$）处处\omterm{def:b1:derivative:def}{可导}，
且 $f'(0) = 0$，但 $f'(x) = 2x
\sin\frac1x - \cos\frac 1x$ 在 $0$ 处没有极限
（见\cref{exo:b1:derivative:2}）。
\end{example}

\begin{example}[恰在一点可导]\label{ex:b1:derivative:onepoint}
设 $f(x) = x^2$（当 $x \in \Q$ 时）且 $f(x) = 0$（当 $x \notin
\Q$ 时）。在 $0$ 处，$\bigl|\frac{f(h) - 0}{h}\bigr| \leq \abs h \to 0$，
故 $f$ 在 $0$ 处\omterm{def:b1:derivative:def}{可导}，且 $f'(0) = 0$。在任一 $x_0
\neq 0$ 处，$f$ 甚至不\omterm{def:b1:continuity:continuous}{连续}：趋于 $x_0$ 的有理数列与无理数列
分别把 $f$ 送到 $x_0^2 \neq 0$ 与
$0$（由\omterm{def:b1:topology:dense}{稠密}性，见\cref{thm:b1:reals:density}）。因此\omterm{def:b1:derivative:def}{可导}性确实是
一个\emph{逐点}的概念：它可以
恰在 $\R$ 中的一点成立，而在其他任何点都不成立。由此得到的实践原则是：
单调性判别法、罗尔定理等结论，
要求导数\emph{在\omterm{prop:b1:reals:intervals}{区间}上}存在；
仅在孤立点处有 $f'(x_0)$，无论有多少个，
都推不出任何整体性结论。
\end{example}

\begin{theorem}[运算法则]\label{thm:b1:derivative:operations}
设 $f, g$ 在 $x_0$ 处\omterm{def:b1:derivative:def}{可导}（公式有意义时），
则有：
\[
(f + g)' = f' + g', \qquad
(fg)' = f'g + fg', \qquad
\Bigl(\frac fg\Bigr)' = \frac{f'g - fg'}{g^2},
\]
又若 $g$ 在 $f(x_0)$ 处\omterm{def:b1:derivative:def}{可导}：
$\;(g \circ f)'(x_0) = g'\bigl(f(x_0)\bigr)\, f'(x_0)$（链式法则
\index{链式法则}）。
\end{theorem}

\begin{proof}
和：显然。乘积：写出
\[
f(x)g(x) - f(x_0)g(x_0)
= \bigl(f(x) - f(x_0)\bigr) g(x) + f(x_0)\bigl(g(x) - g(x_0)\bigr),
\]
两边除以 $x - x_0$，并令 $x \to x_0$（$g$ 在 $x_0$ 处\omterm{def:b1:continuity:continuous}{连续}）。
商的情形：对 $\frac 1g$，有 $\frac{1/g(x) - 1/g(x_0)}{x - x_0} =
\frac{-1}{g(x)g(x_0)}\cdot\frac{g(x) - g(x_0)}{x - x_0}$，再用乘积法则。
链式法则：记 $y_0 = f(x_0)$，定义
$\theta(y) = \frac{g(y) - g(y_0)}{y - y_0}$（当 $y \neq y_0$ 时），
$\theta(y_0) = g'(y_0)$；$\theta$ 在 $y_0$ 处\omterm{def:b1:continuity:continuous}{连续}，并对
$x \neq x_0$，有
\[
\frac{g(f(x)) - g(f(x_0))}{x - x_0}
= \theta\bigl(f(x)\bigr)\cdot \frac{f(x) - f(x_0)}{x - x_0}
\longrightarrow g'(y_0)\, f'(x_0),
\]
其中第一个因子由极限的复合得到（这一手法能干净地处理
情形 $f(x) = f(x_0)$，而朴素的“乘以再除以
$f(x) - f(x_0)$”的做法便会失效）。
\end{proof}

\begin{theorem}[反函数的导数]\label{thm:b1:derivative:inverse}
设 $f$ 在 $I$ 上\omterm{def:b1:continuity:continuous}{连续}且严格单调，并在 $x_0$ 处\omterm{def:b1:derivative:def}{可导}，
且 $f'(x_0) \neq 0$。则 $f^{-1}$
（见\cref{thm:b1:continuity:bijection}）在 $y_0 =
f(x_0)$ 处\omterm{def:b1:derivative:def}{可导}，且
\[
(f^{-1})'(y_0) = \frac{1}{f'(x_0)} = \frac{1}{f'\bigl(f^{-1}(y_0)\bigr)} .
\]
若 $f'(x_0) = 0$，则逆函数在 $y_0$ 处有铅直切线。
\end{theorem}

\begin{proof}
当 $y \to y_0$ 时，记 $x = f^{-1}(y)$：$f^{-1}$ 的\omterm{def:b1:continuity:continuous}{连续}性给出 $x
\to x_0$，且
\[
\frac{f^{-1}(y) - f^{-1}(y_0)}{y - y_0}
= \frac{x - x_0}{f(x) - f(x_0)}
= \frac{1}{\dfrac{f(x) - f(x_0)}{x - x_0}}
\longrightarrow \frac{1}{f'(x_0)} .
\]
铅直切线的断言：若 $f'(x_0) = 0$，则上式中的
差商是一个趋于 $0$ 并保持固定符号的量的倒数
（若 $f$ 严格递增，则 $\frac{f(x) -
f(x_0)}{x - x_0} > 0$ 对一切 $x \neq x_0$ 成立）：因此
$f^{-1}$ 的差商趋于 $+\infty$（趋于 $-\infty$，
对应 $f$ 递减的情形）。逆函数仍\omterm{def:b1:continuity:continuous}{连续}，但
在 $y_0$ 处不\omterm{def:b1:derivative:def}{可导}——其图像是 $f$ 关于对角线的反射，
在 $f$ 水平之处恰好竖直，
如 $x^{1/3}$ 在 $0$ 处与 $x^3$ 的对照所示。
\end{proof}

\begin{example}[两次求逆函数的导数]\label{ex:b1:derivative:inverseaction}
该定理无须求极限，即可重新算出经典函数的导数。
取 $\ln = \exp^{-1}$：在 $y = \eu^x$ 处，
\[
(\ln)'(y) = \frac{1}{\exp'(x)} = \frac{1}{\eu^{x}} = \frac1y ,
\]
对一切 $y > 0$ 均成立，因 $\exp' = \exp$ 永不为零。
取 $\arctan = \tan^{-1}$：在 $y = \tan x$ 处，
\[
(\arctan)'(y) = \frac{1}{1 + \tan^2 x} = \frac{1}{1 + y^2} ,
\]
利用 $\tan' = 1 + \tan^2 > 0$。本例的要点在于：此公式
将关于函数的知识，
换成关于其逆函数的知识，代价只是一次代换——
而该代换（$x = \ln y$，$x = \arctan y$）
恰好体现两个变量分处
同一\omterm{def:b1:logic:inj}{双射}的两侧。
\end{example}

\section{罗尔定理与中值定理}

\begin{proposition}[内点极值]\label{prop:b1:derivative:fermat}
若 $f$ 在 $x_0$ 处\omterm{def:b1:derivative:def}{可导}，且该点为 $I$ 的\emph{内点}，
在该点取得局部极值，则 $f'(x_0) = 0$。
\end{proposition}

\begin{proof}
不妨设取得局部极大值：存在 $r > 0$，使得 $f(x) \leq f(x_0)$
当 $\abs{x - x_0} \leq r$ 时成立，且由内点性，
$x_0$ 两侧的点均属于 $I$。对 $0 < h \leq r$，
差商 $\frac{f(x_0 + h) - f(x_0)}{h}$ 的分子 $\leq 0$，
分母 $> 0$，故差商 $\leq 0$，其极限 $f'(x_0)$
也满足 $\leq 0$（非严格不等式在取极限后保持，
见\cref{thm:b1:seq:order}）；对 $-r \leq h < 0$，差商
$\geq 0$，从而 $f'(x_0) \geq 0$。因此 $f'(x_0) = 0$。（在
端点处只有一侧可用，结论在端点处不成立——
例如 $x$ 在 $\intcc{0}{1}$ 上，在 $1$ 处取最大值，导数为
$1$。）
\end{proof}

\begin{theorem}[罗尔定理]\label{thm:b1:derivative:rolle}
设 $f$ 在 $\intcc{a}{b}$ 上\omterm{def:b1:continuity:continuous}{连续}，
在 $\intoo{a}{b}$ 内\omterm{def:b1:derivative:def}{可导}，且 $f(a) = f(b)$。则必有 $f'(c) = 0$，其中 $c \in
\intoo{a}{b}$。\index{罗尔定理}
\end{theorem}

\begin{proof}
由最值定理（见\cref{thm:b1:continuity:evt}），$f$
在 $\intcc{a}{b}$ 上取得最大值与最小值。若二者
均在端点处取得，则（由于 $f(a) = f(b)$）最大值 $=$ 最小值，$f$
为常数：任一内点 $c$ 均可。否则极值
在某内点 $c$ 处取得，
由\cref{prop:b1:derivative:fermat}可得 $f'(c) = 0$。
\end{proof}

\begin{theorem}[中值定理]\label{thm:b1:derivative:mvt}
设 $f$ 在 $\intcc{a}{b}$ 上\omterm{def:b1:continuity:continuous}{连续}，
在 $\intoo{a}{b}$ 内\omterm{def:b1:derivative:def}{可导}。则存在 $c \in \intoo{a}{b}$，使得
\[
f(b) - f(a) = f'(c)\,(b - a) .
\]
\emph{中值不等式：}若进一步 $m \leq f' \leq M$ 在
$\intoo{a}{b}$ 上成立，则 $m(b-a) \leq f(b) - f(a) \leq M(b-a)$；
特别地，若 $\abs{f'} \leq K$，则 $f$ 为 $K$-利普希茨函数。
\index{中值定理}
\end{theorem}

\begin{proof}
对 $g(x) = f(x) - \frac{f(b) - f(a)}{b - a}(x - a)$ 应用罗尔定理：$g$
在 $\intcc{a}{b}$ 上\omterm{def:b1:continuity:continuous}{连续}，在\omterm{def:b1:topology:closure}{内部}\omterm{def:b1:derivative:def}{可导}，且 $g(a) =
f(a) = g(b)$。取点 $c$，使 $g'(c) = 0$，则 $f'(c) =
\frac{f(b)-f(a)}{b-a}$。由 $f'(c)$ 估界即得该不等式；
利普希茨性质则对任意两点应用此结论。
\end{proof}

\begin{omfigure}
\begin{tikzpicture}
\begin{axis}[omaxis, width=8.2cm, height=5.4cm,
    xmin=0, xmax=4.6, ymin=0, ymax=3.2,
    xtick={0.7,2.2,3.9}, xticklabels={$a$,$c$,$b$}, ytick=\empty]
  \addplot[omDef, very thick, domain=0.7:3.9, samples=100]
    {0.5 + 1.4*ln(x+0.3)};
  \addplot[omExo, thick] coordinates {(0.7,1.5) (3.9,2.5)};
  \addplot[omThm, thick, dashed, domain=1.1:3.4]
    {2.06 + 0.3125*(x - 2.2)};
  \fill[omThm] (axis cs:2.2,2.06) circle (1.6pt);
\end{axis}
\end{tikzpicture}

{\small 中值定理：某条切线（虚线）与
弦（灰色）平行。其横坐标 $c$ 是将罗尔定理
用于函数减去弦之后所得的临界点。}
\end{omfigure}

\begin{example}[牛顿法即海伦迭代]\label{ex:b1:derivative:newton}
求解 $f(x) = 0$ 的牛顿法，是以曲线在当前猜测值
$x_n$ 处的切线代替曲线，并取该切线的
根作为下一个猜测值：
\[
0 = f(x_n) + f'(x_n)(x_{n+1} - x_n)
\quad\Longrightarrow\quad
x_{n+1} = x_n - \frac{f(x_n)}{f'(x_n)} .
\]
取 $f(x) = x^2 - 2$ 运行该方法：
\[
x_{n+1} = x_n - \frac{x_n^2 - 2}{2x_n}
= \frac{x_n}{2} + \frac{1}{x_n}
= \frac12\Bigl(x_n + \frac{2}{x_n}\Bigr) :
\]
恰好就是海伦迭代（见\cref{ex:b1:seq:heron}），
早了两千年。此前观察到的二次收敛速度，如今
可由切线图像解释：在单根附近，曲线与切线
之差为二阶小量，因此每一步大致将误差
平方——一般性结论可由
\cref{ch:b1:taylor}中的泰勒估界推出。本例的要点在于：
二分法（见\cref{ex:b1:continuity:dichotomy}）仅用到\omterm{def:b1:continuity:continuous}{连续}性，每步
只增加一位二进制精度，而牛顿法则付出一阶导数的代价，
使每步的正确数位翻倍。
\end{example}

\begin{example}[作为估计工具的中值定理]\label{ex:b1:derivative:mvtestimate}
$\sqrt{101}$ 有多大？对 $f(t) = \sqrt t$
在 $\intcc{100}{101}$ 上应用中值定理，存在 $c \in \intoo{100}{101}$，
\[
\sqrt{101} - 10 = \frac{1}{2\sqrt c},
\qquad\text{于是}\qquad
\frac{1}{2\sqrt{101}} < \sqrt{101} - 10 < \frac{1}{20} = 0.05 ,
\]
又由 $\sqrt{101} < 10.05$，左\omterm{def:b1:topology:closure}{边界}大于
$\frac{1}{20.1} > 0.0497$，故 $10.0497 < \sqrt{101} < 10.05$
（真值为 $10.049875\dots$）——仅一次导数求值，即得三位
正确小数。同理，$\abs{\sin a - \sin b} \leq
\abs{a - b}$（取 $\abs{\cos}\leq 1$ 作界）：
自\cref{ch:b1:seq}以来所用的 利普希茨 估计，全都是本定理的应用。
要点：中值定理是零阶泰勒公式
——它把一个未知点 $c$ 换成一个不易处理的不等式，
而\cref{ch:b1:taylor}将恰好反复进行这一交换。
\end{example}

\begin{corollary}[单调性判别法]\label{cor:b1:derivative:monotone}
设 $f$ 在 $I$ 上\omterm{def:b1:continuity:continuous}{连续}，在其\omterm{def:b1:topology:closure}{内部}\omterm{def:b1:derivative:def}{可导}。
\begin{enumerate}
  \item 在\omterm{def:b1:topology:closure}{内部} $f' \geq 0$ $\iff$ $f$ 单调递增；$f' = 0$
        $\iff$ $f$ 为常值函数。
  \item 若 $f' > 0$，除有限个使其为零的点外，
        $f$ 是\emph{严格}单调递增的。
\end{enumerate}
\end{corollary}

\begin{proof}
若 $f' \geq 0$：对 $x < y$ 且二者均在 $I$ 中，由中值定理在
$\intcc{x}{y}$ 上得 $f(y) - f(x) = f'(c)(y - x) \geq 0$。反之，
单调递增函数的差商都 $\geq 0$，故其
极限亦然。常值情形：对 $f'$ 与
$-f' \geq 0$ 应用前一结论。严格情形：$f$ 单调递增；若 $f(x) =
f(y)$ 对 $x < y$ 成立时，将使 $f$ 在 $\intcc{x}{y}$ 上恒定，从而迫使 $f' =
0$ 在该\omterm{prop:b1:reals:intervals}{区间}上成立——这将有无穷多个零点。
\end{proof}

\begin{example}[导数相等，函数不等]\label{ex:b1:derivative:intervalmatters}
在 $\R^* = \intoo{-\infty}{0} \cup \intoo{0}{+\infty}$ 上，两个函数
$f(x) = \ln\abs x$ 与 $g(x) = \ln\abs x + \mathbf{1}_{x>0}$
（仅在右半直线上加 $1$）均满足 $f' = g' =
\frac1x$。二者并不相差一个常数：
判别法中的“$f' = 0 \implies f$ 为常数”是关于\emph{\omterm{prop:b1:reals:intervals}{区间}}的结论，
——其证明在两点之间应用中值定理，
而这要求连接两点的整条线段都落在
定义域内。在每条半直线上分别考虑，
$\frac1x$ 的原函数为 $\ln\abs x + c$，每条半直线各一个常数，
共两个独立的常数。\Cref{ch:b1:integration}
承袭了此细节：函数的“原函数”
仅在相差一个常数的意义下确定，且\emph{在其定义域的每个
\omterm{prop:b1:reals:intervals}{区间}上}，而积分表默认地假设了连通性。
\end{example}

\begin{example}[免费的严格性]\label{ex:b1:derivative:strictness}
$x \mapsto x^3$ 在 $\R$ 上是\emph{严格}单调递增的，
尽管其导数在 $0$ 处为零：判别法中的条款
“$f' > 0$ 除有限个点外”正是专为
此类平点而设。相比之下，单独的 $f' \geq 0$ 只给出
宽松意义下的递增（常值函数也满足），并且
若导数在某整个子\omterm{prop:b1:reals:intervals}{区间}上为零，则会使
函数在该处保持不变。实用的规则是：要断言
严格单调性，列出 $f'$ 的零点：有限个（或更一般地，
在任一子\omterm{prop:b1:reals:intervals}{区间}上都无零点）无妨，一整段\omterm{prop:b1:reals:intervals}{区间}上的零点则
是致命的。
\end{example}

\begin{example}[完整的单调性分析]\label{ex:b1:derivative:variations}
研究 $f(x) = x^3 - 3x + 1$ 在 $\R$ 上的情况。求导得 $f'(x) =
3(x^2 - 1)$，在 $\intoo{-\infty}{-1}$ 上为正，在
$\intoo{-1}{1}$ 上为负，在 $\intoo{1}{+\infty}$ 上为正：由
单调性判别法，$f$ 先递增，再递减，再
递增，且有极大值 $f(-1) = 3$ 与极小值
$f(1) = -1$。极限：$\mp\infty$ 对应 $\mp\infty$。由此，
借助单调性表，并在每段单调分支上应用介值定理，
可知 $f$ 恰在以下
各\omterm{prop:b1:reals:intervals}{区间}中各有一个零点：
\[
\intoo{-\infty}{-1}, \qquad \intoo{-1}{1}, \qquad
\intoo{1}{+\infty}
\]
（分界点处函数值异号：$3 > 0 > -1$），
因此方程 $x^3 - 3x + 1 = 0$ 恰有三个实根，
数值上它们约在 $-1.88$、$0.35$、$1.53$ 附近。
要点：单调性表是一种\emph{证明工具}，
而非草图——单调分支加上符号改变，
即恰有一个根，且表格列出全部分支，
穷尽无遗。
\end{example}

\begin{theorem}[莱布尼茨公式]\label{thm:b1:derivative:leibniz}
若 $f, g$ 均 $n$ 阶\omterm{def:b1:derivative:def}{可导}，则 $fg$ 亦然，
且\index{莱布尼茨公式}
\[
(fg)^{(n)} = \sum_{k=0}^{n} \binom nk f^{(k)}\, g^{(n-k)} .
\]
\end{theorem}

\begin{proof}
对 $n$ 用归纳法，与二项式定理完全平行。
$n = 1$ 的情形即乘积法则。设公式对
$n$ 成立，再求一次导数：
\[
(fg)^{(n+1)} = \sum_{k=0}^{n} \binom nk
\Bigl( f^{(k+1)} g^{(n-k)} + f^{(k)} g^{(n-k+1)} \Bigr),
\]
再以 $j = k + 1$ 对第一个和式换标，并归并
$f^{(j)} g^{(n+1-j)}$ 的系数：它是 $\binom{n}{j-1} +
\binom nj = \binom{n+1}{j}$，由帕斯卡法则
（见\cref{prop:b1:counting:identities}），\omterm{def:b1:topology:closure}{边界}项 $j = 0$
与 $j = n + 1$ 处携带的系数 $\binom{n+1}{0} = \binom{n+1}{n+1} =
1$，正如应当的。
\end{proof}

\begin{example}[莱布尼茨公式的应用]\label{ex:b1:derivative:leibnizaction}
计算 $\bigl(x^2 \eu^x\bigr)^{(n)}$，其中 $n \geq 2$。取 $f =
x^2$，其导数很快为零（$f' = 2x$，$f'' = 2$，
$f^{(k)} = 0$，$k \geq 3$），并取 $g = \eu^x$。
莱布尼茨和式中只有三项保留下来，
\[
\bigl(x^2\eu^x\bigr)^{(n)}
= \binom n0 x^2 \eu^x + \binom n1 (2x)\,\eu^x + \binom n2\,
2\,\eu^x
= \eu^x\bigl(x^2 + 2nx + n(n-1)\bigr).
\]
验算：当 $n = 1$ 时，$\eu^x(x^2 + 2x)$，这确实就是
$(x^2\eu^x)'$。小结：一个因子为\omterm{def:b1:poly:def}{多项式}时，宜用莱布尼茨公式，
此时和式只有 $\deg + 1$ 项，
于是所得为闭式，而非抽象的恒等式。（对于两个
求导不会衰减的因子，如 $\eu^x\sin x$，
则更好的工具是\cref{ch:b1:complex}中的复指数。）
\end{example}

\section{凸性}

\begin{definition}\label{def:b1:derivative:convex}
称 $f \colon I \to \R$ 为\emph{凸函数}\index{凸函数}，若
每条弦都位于图像之上：
\[
\forall x, y \in I,\ \forall t \in \intcc{0}{1}, \quad
f\bigl(tx + (1-t)y\bigr) \leq t f(x) + (1-t) f(y).
\]
（$f$ 为\emph{凹函数}，当 $-f$ 为凸函数时。）
\end{definition}

\begin{theorem}[凸性的微分刻画]\label{thm:b1:derivative:convexchar}
设 $f$ 在 $I$ 上\omterm{def:b1:derivative:def}{可导}。下列条件等价：
\begin{enumerate}
  \item $f$ 为\omterm{def:b1:derivative:convex}{凸函数}；
  \item $f'$ 在 $I$ 上单调递增；
  \item 图像位于每条切线之上：$f(y) \geq f(x) +
        f'(x)(y - x)$，对所有 $x, y \in I$。
\end{enumerate}
若 $f$ 二阶\omterm{def:b1:derivative:def}{可导}，则 $f$ 为\omterm{def:b1:derivative:convex}{凸函数} $\iff f'' \geq 0$。
\end{theorem}

\begin{proof}
（1 $\Rightarrow$ 3）凸性可写为 $\frac{f(x + t(y-x)) -
f(x)}{t} \leq f(y) - f(x)$，对 $t \in \intoc{0}{1}$；令 $t \to 0^+$：
$f'(x)(y - x) \leq f(y) - f(x)$。

（3 $\Rightarrow$ 2）设 $x < y$，由 $x$ 与 $y$ 处的两个切线不等式
可得 $f'(x)(y-x) \leq f(y) - f(x) \leq f'(y)(y - x)$，
从而 $f'(x) \leq f'(y)$。

（2 $\Rightarrow$ 1）取定 $x < y$ 及 $t \in \intoo{0}{1}$，并令 $z
= tx + (1-t)y \in \intoo{x}{y}$。在
$\intcc{x}{z}$ 与 $\intcc{z}{y}$ 上由中值定理，存在 $c_1 < z < c_2$，使得
\[
\frac{f(z) - f(x)}{z - x} = f'(c_1) \leq f'(c_2)
= \frac{f(y) - f(z)}{y - z} ,
\]
去分母（利用 $z - x = (1-t)(y-x)$，$y - z = t(y-x)$），
恰好整理为凸性不等式。

二阶\omterm{def:b1:derivative:def}{可导}的情形：$f'' \geq 0 \iff f'$ 单调递增
（见\cref{cor:b1:derivative:monotone}）。
\end{proof}

\begin{omfigure}
\begin{tikzpicture}
\begin{axis}[omaxis, width=8.2cm, height=5.6cm,
    xmin=-2.4, xmax=2.6, ymin=-0.8, ymax=3.4,
    xtick=\empty, ytick=\empty]
  \addplot[omDef, very thick, domain=-2.1:2.3, samples=100]
    {0.5*x*x};
  \addplot[omExo, thick] coordinates {(-1.6,1.28) (2,2)};
  \addplot[omThm, thick, dashed, domain=-1.2:2.4] {0.72*(x-0.72)+0.2592};
\end{axis}
\end{tikzpicture}

{\small 凸性的两个方面：每条弦（灰色）都位于图像之上，
图像位于每条切线（虚线）之上。}
\end{omfigure}

\begin{example}[经典的凸性不等式]\label{ex:b1:derivative:convexineq}
$\exp$ 为\omterm{def:b1:derivative:convex}{凸函数}（$\exp'' = \exp > 0$）：它在 $0$ 处的切线给出
$\eu^x \geq 1 + x$，对所有 $x$ 成立。$\ln$ 为凹函数：它在 $1$ 处的切线
给出 $\ln x \leq x - 1$；其弦则对 $0 < a \leq b$，给出
几何平均与算术平均之间的不等式：在凹性中取 $t =
\frac12$，
\[
\ln\frac{a + b}{2} \geq \frac{\ln a + \ln b}{2} = \ln\sqrt{ab},
\qquad\text{从而}\qquad
\sqrt{ab} \leq \frac{a+b}{2} .
\]
一般的算术--几何平均不等式为
\cref{exo:b1:derivative:9}。
\end{example}

\begin{example}[从零开始的凸性不等式]\label{ex:b1:derivative:xlnx}
函数 $f(t) = t\ln t$ 在 $\intoo{0}{+\infty}$ 上为\omterm{def:b1:derivative:convex}{凸函数}：
$f''(t) = \frac1t > 0$。其中点不等式乘以
$2$，即对所有 $a, b > 0$，
\[
a\ln a + b\ln b \;\geq\; (a + b)\,\ln\frac{a + b}{2} ,
\]
当且仅当 $a = b$ 时等号成立（严格凸性）。试算：$a =
1$，$b = 3$ 得 $3\ln 3 = 3.296$，对比 $4\ln 2 = 2.773$。
这一平凡的不等式，是
\emph{熵}比较的两点情形，它与
琴生不等式（见\cref{exo:b1:derivative:9}）一同再现，并出现在
本科三年级分册中的信息论渐近理论。小结：
要构造一个不等式，需找一个函数，其二阶导数
具有确定符号，并写下凸性所断言的——
内容：微分刻画把一次符号检验化为
无穷多个不等式。
\end{example}

\begin{remark}[导数的常见误区]
(i) \emph{一点处导数为正，
并不能推出其附近的单调性}：$f(x) = \frac x2 + x^2\sin\frac1x$（其中
$f(0) = 0$）时有 $f'(0) = \frac12 > 0$，但
\[
f'(x) = \frac12 + 2x\sin\frac1x - \cos\frac1x
\]
其值为 $-\frac12$，对每个 $x_n = \frac{1}{2\pi n}$ 均成立：
因此 $0$ 的每个\omterm{def:b1:topology:open}{邻域}内都含有下降点。单调性需要 $f'
\geq 0$ \emph{在\omterm{prop:b1:reals:intervals}{区间}上}
（\cref{cor:b1:derivative:monotone}）；逐点符号
只能控制切线的穿越情形。(ii) \emph{罗尔定理的
三个假设全部起作用}：$\abs x$ 在 $\intcc{-1}{1}$ 上
（\omterm{def:b1:topology:closure}{内部}不\omterm{def:b1:derivative:def}{可导}），$x$ 在 $\intcc{0}{1}$ 上（端点处函数值不等），
以及 $x - \lfloor x\rfloor$ 在 $\intcc{0}{1}$ 上（\omterm{def:b1:continuity:continuous}{连续}性
在 $1$ 处失效），三者各恰好破坏一个假设，并使
结论不成立。(iii) \emph{导数可以不\omterm{def:b1:continuity:continuous}{连续}，但
并非任意不\omterm{def:b1:continuity:continuous}{连续}}：$f'$ 可以振荡
（\cref{ex:b1:derivative:notc1}），却总满足
介值性质（达布，
见\cref{exo:b1:derivative:10}）：导数绝不跳跃——若
计算带跳跃的单侧“导数极限”，
则相当于对不\omterm{def:b1:derivative:def}{可导}函数求了导。(iv) \emph{反函数
求导公式需要 $f' \neq 0$}：$x \mapsto x^3$ 是光滑的
严格递增\omterm{def:b1:logic:inj}{双射}，其逆 $x^{1/3}$ 在 $0$ 处有
铅直切线——反函数的\omterm{def:b1:derivative:def}{可导}性恰在
$f'$ 为零处丧失
（见\cref{thm:b1:derivative:inverse}）。
\end{remark}

\begin{remark}[中值定理接下来在哪里起作用]
后续各章中几乎所有的定量结论，
都是本章中值定理的“变装版”：微积分基本定理
（\cref{ch:b1:integration}）对面积函数求导，
并以单调性判别法作结；泰勒–拉格朗日公式
（\cref{ch:b1:taylor}）是中值定理迭代 $n$ 次的结果；
牛顿法的误差分析
与不动点迭代的误差分析（\cref{exo:b1:derivative:11}）
是利普希茨形式；本章的周末问题
（\cref{pb:b1:derivative:1}）把同一利普希茨界转化为
数论问题——\omterm{pb:b1:logic:1}{代数数}与有理数之间的一个排斥不等式，
由此得出历史上第一个\omterm{pb:b1:logic:1}{超越数}。
在本科二年级分册中，中值
\emph{不等式}在多元情形下仍然成立，而等式
则不然。
\end{remark}

\begin{example}[由凹性得到的杨氏不等式]\label{ex:b1:derivative:young}
设 $p, q > 1$，且 $\frac1p + \frac1q = 1$。对所有 $a, b > 0$：
\[
ab \;\leq\; \frac{a^p}{p} + \frac{b^q}{q} .
\]
证明：对 $\ln$ 的凹性应用一次，权重取为
$\frac1p, \frac1q$（即两点形式的琴生不等式，见
\cref{exo:b1:derivative:9}）：
\[
\ln\Bigl(\frac{a^p}{p} + \frac{b^q}{q}\Bigr)
\;\geq\; \frac1p \ln(a^p) + \frac1q \ln(b^q)
= \ln a + \ln b = \ln(ab),
\]
再由 $\ln$ 单调递增，将对数不等式化为
所要证明的结论；等号成立当且仅当 $a^p = b^q$（严格凹性）。当
$p = q = 2$ 时，本式即为算术-几何平均不等式 $ab \leq
\frac{a^2 + b^2}{2}$ 的伪装形式。小结：杨氏
不等式是本科二年级分册中赫尔德与闵可夫斯基
不等式的代数源头——一条凹性陈述，
关于 $\ln$，却被移用到范数上。
\end{example}

\begin{remark}[本册内的展望]
导数在本册结束前还将有三种新的用法。
在\cref{ch:b1:taylor}中，导数被迭代：一点处的 $n$ 次求导
压缩为一个\omterm{def:b1:poly:def}{多项式}加上可控误差，而
中值定理化为拉格朗日余项。在
\cref{ch:b1:curves}中，求导转向几何：对
参数曲线 $t \mapsto (x(t), y(t))$，数对 $(x'(t),
y'(t))$ 是速度\emph{向量}，相切化为
共线，临界点变为待分类的尖点。在
\cref{ch:b1:multivar}中，每次只固定一个变量：偏导数
将本章内容重复两遍，而切线
长成切平面。三章都继承同一套语法：
局部线性逼近加误差项——
这一语法在此首次出现。
\end{remark}

\section{练习}

\begin{exercise}[$\star$]\label{exo:b1:derivative:1}
求下列函数的导数（并指明定义域）：$x^x$；
$\;\ln\bigl(x + \sqrt{x^2+1}\bigr)$; $\;\arctan\frac{1}{x}$;
$\;\sqrt{1 + \eu^{2x}}$.
\end{exercise}

\begin{exercise}[$\star$]\label{exo:b1:derivative:2}
承接\cref{ex:b1:derivative:notc1}，证明：$f(x) = x^2
\sin\frac1x$，$f(0) = 0$，在 $0$ 处\omterm{def:b1:derivative:def}{可导}，且 $f'(0) = 0$，
以及 $f'$ 在 $0$ 处没有极限。
\end{exercise}

\begin{exercise}[$\star$]\label{exo:b1:derivative:3}
利用中值定理或切线不等式，证明
对一切 $x > 0$，成立：
\[
\frac{x}{1 + x} < \ln(1 + x) < x .
\]
由此推出 $\lim_{n\to\infty} \bigl(1 + \frac xn\bigr)^n = \eu^x$ 对
每个 $x > 0$ 成立。
\end{exercise}

\begin{exercise}[$\star$]\label{exo:b1:derivative:4}
设 $P$ 为具有 $k$ 个互异实根的实系数\omterm{def:b1:poly:def}{多项式}。证明
$P'$ 至少有 $k - 1$ 个互异实根，并与
$P$ 的根相互交错。由此推出：若 $P$ 的根全为实数，则 $P'$ 的根也全为实数。
\end{exercise}

\begin{exercise}[$\star\star$]\label{exo:b1:derivative:5}
设 $f$ 在 $\R$ 上\omterm{def:b1:derivative:def}{可导}，且 $f' (x)\to \ell$，其中 $x \to
+\infty$ 时。证明 $\frac{f(x)}{x} \to \ell$ \emph{（用
$\intcc{A}{x}$ 上的中值定理）}。问 $f(x+1) - f(x) \to \ell$ 是否也成立？
\end{exercise}

\begin{exercise}[$\star\star$]\label{exo:b1:derivative:6}
（离散罗尔定理）设 $f$ 为 $n$ 阶\omterm{def:b1:derivative:def}{可导}函数，定义在 $I$ 上，
且在 $n + 1$ 个互异点处取零值。证明 $f^{(n)}$
至少有一个零点。应用：次数 $\leq n$ 的\omterm{def:b1:poly:def}{多项式}若在
$n+1$ 个点处为零，则它恒为零（再次）。
\end{exercise}

\begin{exercise}[$\star\star$]\label{exo:b1:derivative:7}
设 $f$ 在 $\intcc{a}{b}$ 上二阶\omterm{def:b1:derivative:def}{可导}，且 $f(a) = f(b)
= 0$，且 $f(x_0) > 0$，其中 $x_0$ 为某个内点。证明 $f''(c) <
0$ 对某个 $c \in \intoo{a}{b}$ 成立。\emph{（两次应用中值定理，
并比较斜率。）}
\end{exercise}

\begin{exercise}[$\star\star$]\label{exo:b1:derivative:8}
研究函数 $f(x) = \dfrac{\ln x}{x}$ 在
$\intoo{0}{+\infty}$ 上的单调性、极限与最大值。由此推出 $a^b
> b^a$ 对所有满足 $\eu \leq a < b$ 的实数成立，并解决一个著名的特例：
$\eu^\pi$ 与 $\pi^\eu$ 哪个更大？再用较小的整数对
$(2,3)$ 与 $(2,4)$ 加以检验，
并说明它们为何表现不同。
\end{exercise}

\begin{exercise}[$\star\star$]\label{exo:b1:derivative:9}
（算术--几何平均不等式）利用 $\ln$ 的凹性，取一般权重，
（对 $n$ 个点的琴生不等式，需用对 $n$ 的归纳法证明），
证明对正实数 $a_1, \dots, a_n$，有
\[
\sqrt[n]{a_1 a_2 \cdots a_n} \leq \frac{a_1 + \dots + a_n}{n},
\]
当且仅当所有 $a_i$ 相等时取等号。
\end{exercise}

\begin{exercise}[$\star\star\star$]\label{exo:b1:derivative:10}
（达布定理：导数具有介值性）设 $f$ 在 $I$ 上\omterm{def:b1:derivative:def}{可导}，
且 $a < b$ 为 $I$ 中的数，满足 $f'(a) < v < f'(b)$。
考虑 $g(x) = f(x) - vx$，并考虑 $g$ 在
$\intcc{a}{b}$ 上取最小值的点，证明 $f'(c) = v$ 对某个 $c
\in \intoo{a}{b}$ 成立——尽管 $f'$ 未必\omterm{def:b1:continuity:continuous}{连续}
（见\cref{exo:b1:derivative:2}）。
\end{exercise}

\begin{exercise}[$\star\star\star$]\label{exo:b1:derivative:11}
设 $f \colon \R \to \R$ \omterm{def:b1:derivative:def}{可导}，且 $\abs{f'(x)} \leq k
< 1$ 对所有 $x$ 成立（即\emph{压缩\omterm{def:b1:logic:map}{映射}}）。证明 $f$ 恰有一个
不动点 $\ell$，且每个满足 $u_{n+1} = f(u_n)$ 的数列
都收敛到 $\ell$，并满足 $\abs{u_n - \ell} \leq k^n\abs{u_0 -
\ell}$。\emph{（存在性：对
$f(x) - x$ 在一个大\omterm{prop:b1:reals:intervals}{区间}上应用介值定理，并利用利普希茨界；
或利用完备性与柯西准则。）}
\end{exercise}

\begin{exercise}[$\star\star\star$]\label{exo:b1:derivative:12}
（柯西中值定理与洛必达法则）
\begin{enumerate}
  \item 设 $f, g$ 在 $\intcc{a}{b}$ 上\omterm{def:b1:continuity:continuous}{连续}，
        在 $\intoo{a}{b}$ 内\omterm{def:b1:derivative:def}{可导}，且 $g'$ 在其中处处不为零。
        证明 $g(b) \neq g(a)$，并证明存在 $c \in
        \intoo{a}{b}$ 使得下式成立：
        \[
        \frac{f(b) - f(a)}{g(b) - g(a)} = \frac{f'(c)}{g'(c)}
        \]
        \emph{（对 $h = f - \lambda g$ 应用罗尔定理，取适当的
        常数 $\lambda$）}。
  \item 由此推出在某点处 $\frac00$ 型的洛必达法则：
        若 $f(a) = g(a) = 0$，且
        $\frac{f'(x)}{g'(x)} \to \ell$，其中 $x \to a^+$，则
        $\frac{f(x)}{g(x)} \to \ell$。
  \item 说明逆命题不成立：对 $f(x) = x^2\sin\frac1x$
        （$f(0) = 0$）及 $g(x) = x$，商
        $\frac{f}{g}$ 在 $0$ 处有极限，但
        $\frac{f'}{g'}$ 却没有极限。
\end{enumerate}
\end{exercise}

\section{问题：刘维尔不等式与
第一个超越数}

\begin{problem}[{周末问题——\omterm{pb:b1:logic:1}{代数数}排斥有理数：
$\abs{x - p/q} \geq C/q^d$，以及
$\sum 10^{-n!}$ 的超越性}]
\label{pb:b1:derivative:1}
若一个实数是某个非零的
整系数\omterm{def:b1:poly:def}{多项式}的根，则称为\emph{\omterm{pb:b1:logic:1}{代数数}}，否则称为\emph{\omterm{pb:b1:logic:1}{超越数}}。
1844 年，刘维尔构造出有史以来第一个
被\emph{证明}为\omterm{pb:b1:logic:1}{超越数}的数，其证明的核心是本章的中值定理：
次数为 $d$ 的\omterm{pb:b1:logic:1}{代数数}
不能被有理数以优于 $C/q^d$ 的精度逼近——因此一个
能以\emph{快于任何幂次}的速度被逼近的数，不可能是
\omterm{pb:b1:logic:1}{代数数}。本题建立该不等式，构造
刘维尔数 $L = 0.110001000\dots$（其中数字 1 位于阶乘位置，
借助\cref{pb:b1:reals:1}中的数位方法），
证明其超越性，并以康托尔的另一证法作结，
给出 $\sqrt2$ 与
$2^{1/3}$ 的有效界。\index{超越数}\index{刘维尔
不等式}\index{代数数}

\medskip
\noindent\textbf{第一部分——有理数可被
逼近到何种程度？}
\begin{enumerate}
  \item 证明：两个不同的有理数 $\frac ab \neq \frac pq$
        （其中 $b, q \geq 1$）满足 $\bigl|\frac ab -
        \frac pq\bigr| \geq \frac{1}{bq}$。由此推出：若 $x =
        \frac ab$，且 $0 < \bigl|x - \frac pq\bigr| <
        \frac{1}{bq}$，则不存在这样的 $\frac pq$——有理数
        在尺度 $\frac 1q$ 上排斥所有其他有理数。
  \item 证明：对\emph{每一个}有理数 $\frac pq$（$q \geq
        1$）：$\bigl|\sqrt2 - \frac pq\bigr| \geq
        \frac{1}{4q^2}$ \emph{（若距离超过 $1$，结论显然；
        否则估计 $\abs{\sqrt2 + p/q} < 4$，
        并利用非零整数 $\abs{p^2 - 2q^2} \geq 1$）}。
  \item 反过来：验证 $(p, q) \mapsto (p +
        2q, p + q)$ 保持 $\abs{p^2 - 2q^2} = 1$，并由 $(1,1)$ 生成
        数对 $(3,2)$、$(7,5)$、$(17,12)$，
        $(41,29)$、$(99,70)$，并证明每一对都满足
        \[
        \Bigl|\sqrt2 - \frac pq\Bigr| =
        \frac{1}{q^2\,(\sqrt2 + p/q)} < \frac{1}{2q^2} :
        \]
        由此得到无穷多个 $2$ 阶逼近。结合第 2 问，
        可知 $\sqrt 2$ 的逼近指数
        \emph{恰好}为 $2$。
  \item （狄利克雷）设 $x$ 为无理数，$N \in \N^*$。
        考虑 $N + 1$ 个数 $0, x, 2x,
        \dots, Nx$ 的小数部分，落入 $N$ 个小\omterm{prop:b1:reals:intervals}{区间} $\intco{\frac kN}{\frac{k +
        1}{N}}$，由鸽巢原理
        （见\cref{cor:b1:counting:pigeonhole}），必有两个落入同一\omterm{prop:b1:reals:intervals}{区间}。
        由此推出 $q \leq N$ 及 $p$，使得 $\abs{qx - p} <
        \frac 1N$，因此有无穷多个有理数满足
        $\bigl|x - \frac pq\bigr| < \frac{1}{q^2}$：\emph{每个}
        无理数都可被 $2$ 阶逼近。
\end{enumerate}

\medskip
\noindent\textbf{第二部分——刘维尔不等式。}设 $x$ 为
无理数且为\omterm{pb:b1:logic:1}{代数数}。
\begin{enumerate}[resume]
  \item 证明：在以 $x$ 为根的非零整系数\omterm{def:b1:poly:def}{多项式}中，
        存在一个\omterm{def:b1:poly:def}{多项式} $P$，次数为 $d$，
        其\emph{无有理根}；并检验 $d \geq 2$
        \emph{（分离出因式 $X - \frac ab$，其系数属于 $\Q$，并去分母；
        若次数为 $1$，则 $x$ 将
        为有理数）}。
  \item 证明：对每个有理数 $\frac pq$（$q \geq 1$），有
        $\bigl|P\bigl(\frac pq\bigr)\bigr| \geq \frac{1}{q^d}$
        \emph{（$q^d P(p/q)$ 是非零整数）}。
  \item 记 $M = \max_{\intcc{x-1}{x+1}} \abs{P'}$
        （见\cref{thm:b1:continuity:evt}）。利用
        中值定理，在 $x$ 与 $\frac pq$ 之间证明
        \emph{刘维尔不等式}：取 $C = \min\bigl(1,
        \frac 1M\bigr) > 0$，则
        \[
        \Bigl| x - \frac pq \Bigr| \geq \frac{C}{q^{\,d}}
        \qquad\text{对每个有理数 } \frac pq,\ q \geq 1 .
        \]
  \item 称 $x$ 为\emph{刘维尔数}，当对每个 $n \in
        \N$ 都存在有理数 $\frac pq$，满足 $q \geq 2$ 且
        $0 < \bigl|x - \frac pq\bigr| < q^{-n}$。证明刘维尔数
        是无理数\emph{（第 1 题：取
        $n$，使 $2^{\,n-1} > b$）}。
  \item 证明刘维尔定理：\emph{刘维尔数是
        \omterm{pb:b1:logic:1}{超越数}}\emph{（综合第 7 题与第 8 题：
        不等式 $C < q^{\,d-n}$ 对充分大的 $n$ 不成立）}。
\end{enumerate}

\medskip
\noindent\textbf{第三部分——数 $L$。}
\begin{enumerate}[resume]
  \item 设 $L$ 为十进制数字串（按
        \cref{pb:b1:reals:1}的意义）的值，其中数字 $1$ 位于位置 $n!$
        （$n = 1, 2, 3, \dots$）处，其余位置数字为 $0$，即
        $L = \sup_k t_k$，其中 $t_k =
        \sum_{n=1}^{k} 10^{-n!}$。写出前 $25$
        位 $L$ 的数字。
  \item 证明尾项的夹逼估计：对每个 $k \geq 1$，有
        \[
        10^{-(k+1)!} \;\leq\; L - t_k \;\leq\;
        \frac{10}{9}\,10^{-(k+1)!} \;<\; 2\cdot 10^{-(k+1)!}
        \]
        \emph{（将 $t_k$ 之后的每个部分和
        以有限几何和加以控制）}。
  \item 设 $t_k = \frac{p_k}{q_k}$，其中 $q_k = 10^{k!}$。
        证明 $0 < L - \frac{p_k}{q_k} < \frac{2}{q_k^{\,k+1}}$，
        并由此推出 $L$ 是第 8 题意义下的
        刘维尔数。
  \item 得出结论：$L$ 为\omterm{pb:b1:logic:1}{超越数}——这是历史上第一个显式的
        例子（刘维尔，1844 年）。再直接检验其无理性：其数字
        不是最终周期的（间隔逐渐增大，
        见\cref{pb:b1:reals:1}
        第 20 题）。
  \item 推广：将每个数字 $1$ 换成任意非零数字
        $d_n \in \intint{1}{9}$。证明所得值仍是刘维尔数，并
        借助对角线论证（见\cref{pb:b1:reals:1}
        第 22 题）的论证应用于
        这些数字选取——可知这种形式的\omterm{pb:b1:logic:1}{超越数}
        有不可数多个。
\end{enumerate}

\medskip
\noindent\textbf{第四部分——逼近阶的
层次。}称 $x$ 为\emph{$\mu$ 阶可逼近}，当
        存在常数 $c > 0$，使得有无穷多个有理数满足
        $\bigl|x - \frac pq\bigr| < \frac{c}{q^{\mu}}$。
\begin{enumerate}[resume]
  \item 汇总第一至三部分，得到逼近阶的层次：有理数可按 $1$ 阶逼近，
        且不能更好；$\sqrt 2$ 可按
        $2$ 阶逼近，且不能更好；每个无理数可按至少
        $2$ 阶逼近；次数为 $d$ 的\omterm{pb:b1:logic:1}{代数数}
        不能按超过 $d$ 阶逼近；刘维尔数可按任意阶逼近。给出
        每个结论的依据，并引用相关的题。
  \item 证明 $L + r$ 对每个有理数
        $r = \frac ab$ 都是刘维尔数\emph{（平移逼近分数：
        新的分母为 $b\,q_k$）}。
        由此得出：刘维尔数——从而\omterm{pb:b1:logic:1}{超越数}——
        在 $\R$ 中\omterm{def:b1:topology:dense}{稠密}。
  \item （康托，1874 年）证明\omterm{pb:b1:logic:1}{代数数}全体构成的\omterm{def:b1:logic:sets}{集合}
        是可数的：只有有限个整系数\omterm{def:b1:poly:def}{多项式}，其次数与
        $\abs{\text{系数}}$ 之和不超过 $h$，
        每个至多有 $\deg$ 个根。可数个有限集的
        并仍是可数的。由于任何数列都不能穷尽 $\R$
        （见\cref{pb:b1:reals:1}，
        第 22 题）。故\omterm{pb:b1:logic:1}{超越数}存在——事实上它们构成一个
        不可数集。比较两种证明：
        刘维尔的证明能给出哪些康托证明不能给出的内容？
  \item 直接由第 2 题证明 $\sqrt 2$
        \emph{不是}刘维尔数\emph{（对 $n \geq 3$，
        不等式 $q^{-n} > \frac{1}{4q^2}$ 给出 $q$ 的界；
        于是只剩有限个候选有理数，它们
        都与 $\sqrt2$ 的距离为正）}。推广：
        任何\omterm{pb:b1:logic:1}{代数数}都不是刘维尔数。
\end{enumerate}

\medskip
\noindent\textbf{第五部分——有效常数。}
\begin{enumerate}[resume]
  \item 对佩尔数对 $(99, 70)$，验证 $99^2 - 2\cdot70^2
        = 1$，并计算精确误差
        \[
        \sqrt2 - \frac{99}{70}
        = \frac{-1}{70^2\,\bigl(\sqrt2 + \frac{99}{70}\bigr)},
        \qquad
        \Bigl|\sqrt 2 - \frac{99}{70}\Bigr| \approx 7.2\cdot
        10^{-5} :
        \]
        由一个三位数的分数得到五位正确数字。
  \item 将第二部分用于 $x = 2^{1/3}$、$P = X^3 - 2$：验证 $P$
        无有理根，并估计 $M = \max_{\intcc{x-1}{x+1}}
        3t^2 \leq 3\,(1 + 2^{1/3})^2 < 16$，并得出
        有效不等式
        \[
        \Bigl| 2^{1/3} - \frac pq \Bigr| \geq
        \frac{1}{16\,q^3} \qquad \text{ 对所有 } \frac pq .
        \]
  \item 应用：证明任何逼近 $2^{1/3}$ 的有理数
        误差在 $10^{-6}$ 以内时，其分母须满足 $q \geq 40$。
  \item 证明底数 $10$ 并无实质影响：其二进制类比
        $\sum_{n\geq1} 2^{-n!}$（阶乘位置为 1 的
        二进制数字串的值）也是
        刘维尔数，从而是\omterm{pb:b1:logic:1}{超越数}。
\end{enumerate}

\medskip
\noindent\textbf{第六部分——前沿与综合。}
\begin{enumerate}[resume]
  \item 设 $x^\dagger$ 为小数串的值，该串恰在
        位置 $3^k$（$k \geq 0$）处为 1。证明
        $x^\dagger$ 可被 $3$ 阶逼近，
        由刘维尔不等式推出 $x^\dagger$ 既非有理数，也非二次无理数。说明
        此方法为何在此停滞：
        阶 $3$ 与次数 $\geq 3$ 的代数性相容，
        而填补这一缺口
        （对每个\omterm{pb:b1:logic:1}{代数数}，指数 $> 2$ 即可）
        即罗斯定理，远超本册范围。
  \item 将康托尔的论证定量化：证明次数 $\leq d$ 的\omterm{pb:b1:logic:1}{代数数}，
        若由系数属于 $\intint{-H}{H}$ 的\omterm{def:b1:poly:def}{多项式}给出，则其个数至多为
        $d\,(2H + 1)^{d+1}$。
        （正是这一有限性使第 17 题得以成立。）
  \item 综合，每小题一句话作答：（i）指出刘维尔证明中
        唯一的分析要素（本章的哪一定理，用于何处）；
        （ii）陈述支撑它的张力：
        （整性要求 $\abs{P(p/q)} \geq
        q^{-d}$，光滑性禁止 $\abs{P(p/q)} > M\abs{x -
        p/q}$）；（iii）比较刘维尔与康托尔
        关于\omterm{pb:b1:logic:1}{超越数}存在性的证明；（iv）指出本册在何处
        再次触及同一主题：
        \cref{ch:b1:integration}中的周末问题借同样的
        整性与小量之间的夹逼证明 $\pi$ 为无理数，
        只是以积分代替导数。
\end{enumerate}
\end{problem}
